\section{Analytische Struktur des Optimierungsproblems}
\subsection{Motivation}
In klassischer numerischer Optimierung werden Jacobi-Matrizen häufig durch
finite Differenzen oder automatische Differentiation berechnet. Im AIM-Rahmen
ist dies nicht erforderlich: Das Alignment-Modell ist intrinsisch
differenzierbar.

\subsection{Grundlegende Beobachtung}
\begin{proposition}
Alle Zustandsvariablen entstehen durch Bogenlängenintegration:
\[
\theta(s)=\theta_0+\int_0^s\kappa(\sigma)\,\dd\sigma,
\]
\[
\gamma(s)=\gamma_0+\int_0^s(\cos\theta,\sin\theta)\,\dd\sigma.
\]
\end{proposition}

\subsection{Parametertypen}
\begin{definition}
Primäre Parameterklassen sind Krümmungsparameter \(dK\), Längenparameter
\(dL\) und Übergangsparameter \(dT\).
\end{definition}

\subsection{Analytische Ableitungen}
\begin{proposition}
\[
\frac{\partial\theta(s)}{\partial p}
=\int_0^s\frac{\partial\kappa}{\partial p}\,\dd\sigma.
\]
\end{proposition}
\begin{proposition}
\[
\frac{\partial\gamma(s)}{\partial p}
=\int_0^s
\begin{pmatrix}-\sin\theta\\\cos\theta\end{pmatrix}
\frac{\partial\theta}{\partial p}\,\dd\sigma.
\]
\end{proposition}
Alle Ableitungen reduzieren sich auf Bogenlängenintegrale.

\subsection{Zentrale Erkenntnis}
\begin{proposition}
Das Optimierungsproblem besitzt exakte Jacobi-Matrizen ohne numerische
Approximation.
\end{proposition}
Dies beseitigt finite Differenzen, numerische Instabilität und hohen
Rechenaufwand.

\subsection{Blockstruktur}
Der Parametervektor zerfällt in \((\boldsymbol\kappa,\boldsymbol\phi)\) und
induziert
\[
H=\begin{pmatrix}
H_{\kappa\kappa}&H_{\kappa\phi}\\
H_{\phi\kappa}&H_{\phi\phi}
\end{pmatrix}.
\]

\subsection{Interpretation}
\begin{summary}
Die Struktur wird nicht numerisch aufgeprägt, sondern folgt unmittelbar aus
der krümmungsbasierten Modellierung. Daraus entstehen exakte Ableitungen, dünn
besetzte Matrizen und natürliche Blockstruktur -- ein zentraler Vorteil von AIM.
\end{summary}
